% =============================================================================
%  Chapter 10 — Serendipity Element Families
% =============================================================================

\chapter{Serendipity Element Families}
\label{ch:serendipity}

This chapter documents the implementation of the new serendipity element
families added to \texttt{fall\_n}, their integration into
\texttt{ElementGeometry}, the associated tests, and the benchmark results.

The implementation goal was architectural, not merely additive:
to introduce a new family without changing the overall layering of the codebase.
In particular:

\begin{itemize}[nosep]
  \item the concrete element remains a pure value type;
  \item \texttt{ElementGeometry} continues to provide the type-erased boundary;
  \item quadrature remains a separate compile-time policy;
  \item no new virtual dispatch is introduced in the hot integration path.
\end{itemize}


% ─────────────────────────────────────────────────────────────────────────
\section{Why Serendipity?}
\label{sec:serendipity_why}

For tensor-product cells, the quadratic Lagrange family and the quadratic
serendipity family differ substantially in node count:

\begin{center}
\begin{tabular}{lcc}
  \toprule
  Reference cell & Quadratic Lagrange & Quadratic serendipity \\
  \midrule
  Quadrilateral & 9 nodes  & 8 nodes \\
  Hexahedron    & 27 nodes & 20 nodes \\
  \bottomrule
\end{tabular}
\end{center}

The serendipity family removes face-center and volume-center nodes while
retaining vertex and edge nodes.  This makes the element attractive when the
primary objective is to reduce geometry, interpolation, and assembly cost
without dropping to first order.

Architecturally, serendipity is also a good stress test of the current design.
It requires:

\begin{itemize}[nosep]
  \item a new reference cell type;
  \item a new concrete geometry type;
  \item VTK support for new node counts;
  \item optional Gmsh import support;
  \item no changes in the physical element layer.
\end{itemize}

That is precisely the extensibility story claimed in
\cref{sec:extensibility}; this chapter documents the claim becoming real code.


% ─────────────────────────────────────────────────────────────────────────
\section{Mathematical Structure}
\label{sec:serendipity_math}

\subsection{Tensor-product serendipity on the cube}

For the quadratic serendipity family on the reference hypercube
\(\hat{\Omega} = [-1,1]^d\), nodes are located at:

\begin{itemize}[nosep]
  \item the \(2^d\) vertices;
  \item the \(d\,2^{d-1}\) edge midpoints.
\end{itemize}

Thus, for \(d=3\), the quadratic hexahedron has
\[
2^3 + 3 \cdot 2^2 = 8 + 12 = 20
\]
nodes.

\paragraph{Corner shape functions.}
Let \(\mathbf{r} = (r_1,\dots,r_d)\) be a corner node with \(r_j \in \{-1,1\}\).
The quadratic serendipity shape function is
\[
N_{\mathbf{r}}(\mathbf{x})
=
2^{-d}
\Bigg( \prod_{j=1}^{d}(1+r_j x_j) \Bigg)
\Bigg( 1 + \sum_{j=1}^{d} r_j x_j - d \Bigg).
\]

\paragraph{Edge-midpoint shape functions.}
Let \(\mathbf{r}\) be an edge node, so exactly one component \(r_k = 0\) and all
others are \(\pm 1\).  Then
\[
N_{\mathbf{r}}(\mathbf{x})
=
2^{1-d}(1-x_k^2)\prod_{j \ne k}(1+r_j x_j).
\]

These expressions unify the 1D, 2D, and 3D quadratic serendipity family.
The implementation therefore uses a single generic formula for
\texttt{TopDim} \(=1,2,3\), with compile-time node descriptors and
\texttt{consteval} reference-node tables.

\subsection{Why the tetrahedral ``serendipity'' family is different}

The simplex case requires more care.  For tetrahedra at order \(p=2\), the
standard quadratic Lagrange simplex already uses only:

\begin{itemize}[nosep]
  \item 4 vertices;
  \item 6 edge midpoints;
  \item no face-interior nodes;
  \item no volume-interior nodes.
\end{itemize}

Therefore the practical quadratic tetrahedral family is already
``serendipity-like'' in the sense relevant to this codebase: there are no
face-center or body-center DOFs to remove.

For this reason, the new \texttt{SerendipitySimplexElement} is represented as an
explicit family alias over the existing simplex implementation for
orders \(1\) and \(2\).  This is mathematically honest and avoids inventing an
ad hoc reduced tetrahedron with unclear approximation properties.


% ─────────────────────────────────────────────────────────────────────────
\section{Implementation Structure}
\label{sec:serendipity_impl}

\subsection{Reference cells}

The tensor-product family is implemented in
\texttt{SerendipityCell.hh}.  The file introduces:

\begin{itemize}[nosep]
  \item \texttt{serendipity\_num\_nodes(top\_dim, order)};
  \item \texttt{serendipity\_reference\_nodes<TopDim, Order>()};
  \item \texttt{SerendipityBasis<TopDim, Order>};
  \item \texttt{SerendipityCell<TopDim, Order>}.
\end{itemize}

Reference nodes are computed with \texttt{consteval}, so the full node layout is
available at compile time and can participate in static assertions and
compile-time topology generation.

\subsection{Concrete geometry}

The concrete tensor-product family is implemented in
\texttt{SerendipityElement.hh} as
\[
\texttt{SerendipityElement<Dim, TopDim, Order>},
\]
mirroring the shape of the existing \texttt{SimplexElement} and
\texttt{LagrangeElement} abstractions:

\begin{itemize}[nosep]
  \item it stores node IDs and bound node pointers;
  \item it exposes \texttt{H()}, \texttt{dH\_dx()},
        \texttt{map\_local\_point()}, and \texttt{evaluate\_jacobian()};
  \item it carries a \texttt{ReferenceCell} type with compile-time topology.
\end{itemize}

The new family is therefore a first-class Layer~1 geometry type.

\subsection{ElementGeometry integration}

\texttt{ElementGeometry.hh} was extended in only one place:
its face-topology and face-factory logic now recognises the serendipity family.
No changes were required in the outer wrapper API.

This confirms the architectural objective:
the extension point is the concrete geometry type, not the type-erasure layer.

\subsection{Quadrature choices}

The following quadratures were selected from the rules already present in the
library:

\begin{center}
\begin{tabular}{lll}
  \toprule
  Element & Quadrature & Rationale \\
  \midrule
  HEX20 & \texttt{GaussLegendreCellIntegrator<3,3,3>} & full quadratic tensor-product integration \\
  QUAD8 face & \texttt{GaussLegendreCellIntegrator<3,3>} & compatible face integration \\
  EDGE3 face & \texttt{GaussLegendreCellIntegrator<3>} & compatible edge integration \\
  TET10 serendipity & \texttt{SimplexIntegrator<3,2>} & same requirements as quadratic simplex \\
  \bottomrule
\end{tabular}
\end{center}

The choice is intentionally conservative.  The goal of this first integration is
correctness and compatibility, not aggressive reduced integration.


% ─────────────────────────────────────────────────────────────────────────
\section{VTK and Gmsh Support}
\label{sec:serendipity_io}

\subsection{VTK}

\texttt{VTKCellTraits.hh} now recognises:

\begin{itemize}[nosep]
  \item 8-node quadratic serendipity quadrilaterals
        (\texttt{VTK\_QUADRATIC\_QUAD});
  \item 20-node quadratic serendipity hexahedra
        (\texttt{VTK\_QUADRATIC\_HEXAHEDRON}).
\end{itemize}

The test executable writes four validation files:

\begin{itemize}[nosep]
  \item \texttt{data/output/serendipity\_hex20\_sample\_mesh.vtu}
  \item \texttt{data/output/serendipity\_hex20\_sample\_gauss.vtu}
  \item \texttt{data/output/serendipity\_tet10\_sample\_mesh.vtu}
  \item \texttt{data/output/serendipity\_tet10\_sample\_gauss.vtu}
\end{itemize}

These contain both the cell geometry and the materialised Gauss cloud, together
with a simple scalar field for visual inspection in ParaView.

\subsection{Gmsh}

\texttt{GmshDomainBuilder.hh} was extended with support for:

\begin{itemize}[nosep]
  \item \texttt{HEX\_20} (element type 17);
  \item \texttt{QUA\_8} (element type 16) for boundary faces.
\end{itemize}

The chosen internal node ordering matches the standard corner-plus-edge layout,
so no nontrivial permutation was required for these two element types.


% ─────────────────────────────────────────────────────────────────────────
\section{Validation Tests}
\label{sec:serendipity_tests}

The new test executable \texttt{test\_serendipity.cpp} validates:

\begin{enumerate}[nosep]
  \item static node counts for the family;
  \item reference-node coordinates for HEX20;
  \item partition of unity of the HEX20 basis;
  \item Kronecker interpolation property at all 20 reference nodes;
  \item exact reproduction of a representative quadratic polynomial;
  \item unit-cube volume and mixed-monomial integration through
        \texttt{ElementGeometry};
  \item creation and integration of quadratic serendipity face geometry;
  \item partition of unity and volume for the tetrahedral serendipity family;
  \item VTK cell-type dispatch and actual file generation.
\end{enumerate}

This test set is deliberately front-loaded toward interpolation quality.
Before benchmarking or embedding the element into larger analyses, the basis
itself must satisfy the core finite element identities.


% ─────────────────────────────────────────────────────────────────────────
\section{Benchmark Results}
\label{sec:serendipity_benchmark}

The microbenchmark measures repeated geometric integration of a constant field
through \texttt{ElementGeometry}.  It was executed locally after the full test
suite passed.

\begin{center}
\begin{tabular}{lcc}
  \toprule
  Case & Wall time [s] & Observation \\
  \midrule
  HEX20 serendipity & 6.04 & baseline for 4000 integrations \\
  HEX27 triquadratic & 19.14 & about \(3.2\times\) slower than HEX20 \\
  TET10 serendipity(simplex) & 0.334 & practically same as quadratic simplex \\
  TET10 quadratic simplex & 0.364 & expected baseline \\
  \bottomrule
\end{tabular}
\end{center}

\paragraph{Interpretation.}
The HEX20 result is the important one.  The node reduction from 27 to 20 lowers
the cost of geometry interpolation and Jacobian assembly enough to be visible
even though both elements use the same \(3 \times 3 \times 3\) Gauss rule.

The tetrahedral result confirms the modelling claim made earlier:
at quadratic order, ``serendipity tetrahedron'' and the existing quadratic
simplex are effectively the same finite element family in this codebase.

\paragraph{Iteration after benchmarking.}
No further optimisation was required after the first measurement because the
results already matched the intended architectural outcome:

\begin{itemize}[nosep]
  \item reduced tensor-product node count translated into lower runtime cost;
  \item simplex serendipity remained mathematically and computationally aligned
        with the existing implementation;
  \item no additional runtime indirection was introduced.
\end{itemize}


% ─────────────────────────────────────────────────────────────────────────
\section{Architectural Takeaways}
\label{sec:serendipity_takeaways}

This extension validates several design choices of the current library:

\begin{enumerate}[nosep]
  \item \textbf{External polymorphism works.}
        A new family was added without redesigning \texttt{ElementGeometry}.
  \item \textbf{Reference-cell-centric design scales.}
        Static reference nodes and subentity topology remain the right place to
        encode geometric knowledge.
  \item \textbf{Geometry and physics remain decoupled.}
        The new elements plug into the same generic physical element and VTK
        machinery.
  \item \textbf{The simplex/tensor-product split is conceptually sound.}
        The tetrahedral case did not need an artificial reduced basis, while
        the hexahedral case benefited from a genuinely different family.
\end{enumerate}


% ─────────────────────────────────────────────────────────────────────────
\section{Future Work}
\label{sec:serendipity_future}

The present implementation intentionally stops at order \(2\).  Sensible next
steps would be:

\begin{itemize}[nosep]
  \item reduced-integration variants for HEX20;
  \item serendipity quadrilateral usage in shell formulations;
  \item higher-order incomplete tensor-product families if they are needed;
  \item a more explicit distinction in the type system between
        ``quadratic simplex'' and ``serendipity simplex'' if higher-order
        simplex variants are later introduced.
\end{itemize}

At this stage, however, the library now has a complete and usable quadratic
serendipity pathway for both the tensor-product and simplex branches, with
tests, benchmark coverage, and VTK output.
